% Chapter Template

\section{Introduction} % Main chapter title

\label{ChapterX} % Change X to a consecutive number; for referencing this chapter elsewhere, use \ref{ChapterX}

%----------------------------------------------------------------------------------------
%	SECTION 1
%----------------------------------------------------------------------------------------
We develop a macroeconomic stock flow consistent agent based model (SFC-ABM) of the Dutch economy which we calibrated with the Statistics Netherlands microdata database and Eurostat national accounts data. The CBS database contains administrative data about the income, paid taxes, received and paid social transfers and wealth of every Dutch citizen and household and hence is a valuable source of information for macroeconomic policy recommendations\footnote{We strictly follow the CBS law and the Dutch privacy law (Algemene Verordening Gegevensbescherming (AVG)). The anonymized microdata are immediately aggregated within the strict secure environment of CBS and only the aggregated data are then used in all subsequent steps. The aggregation is sufficient that that no individual data can be extracted from this paper.}\footnote{The Microdata database also has firm data which we can incorporate in the model in the future. Currently, the firm data is calibrated using Eurostat national accounts data.}. To our knowledge current standard modelling approaches, for instance heterogeneous agent dynamic stochastic general equilibrium (DSGE) models or more specifically Heterogeneous Agent New Keynesian (HANK) models like the model of \cite{kaplan_monetary_2018}, are not capable of using this degree of granularity without aggregating the data into higher level moments. With this paper we want to show that agent-based models can complement traditional economic modeling practices. ABMs can relax disputed assumptions such as equilibrium, representative agents, rational expectations, and linearity and could therefore potentially address the shortcomings of macroeconomic models in forecasting economic crises \citep{fagiolo_macroeconomic_2017}. Additonally, the framework is flexible enough to allow for many different assumptions regarding firm and household behaviour and their interactions. We show this by incorporating an optimal consumption model within the agent-based model.\\

\noindent The model builds upon \cite{poledna_economic_2019}. The agent based model incorporates all economic activities and economic entities classified by the European System of Accounts. This means that every legal entity and natural person and every firm is represented by a separate heterogeneous agent. Markets are fully decentralized and economic interactions are modelled by search-and- matching processes. We depart from \cite{poledna_economic_2019} by using a household model derived from \cite{kaplan_monetary_2018}. Household decisions become forward looking and dependent upon income, illiquid capital, liquid capital holdings, the interest rates of illiquid capital, the interest rate of liquid capital, and inflation. To analyse the external validity of the model, we will compare the forecasting capabilities with benchmark models which are an auto-regressive (AR) model, a vector auto-regressive (VAR) model, an auto-regressive model with exogenous predictors (ARX), and an vector auto-regressive model with exogenous predictors (VARX). We will also compare the adjusted agent based model with the agent based model of \cite{poledna_economic_2019} which we call the base agent based model and an agent based model with a more complex consumption heuristic. We also analyze the internal validity of the model by comparing the auto- and cross-correlations of the model generated paths of the macroeconomic variables with their empirical counterparts. \\

\noindent The model is composed of the following elements. Firms are part of an input-output network. Households are boundedly rational and use a simple AR(1) expectation rule which is used to make consumption and savings decisions. We will use a boundedly rational forward looking optimisation problem such that we can model different marginal propensities to consume for different levels of income, and of liquid and illiquid assets. The households are boundedly rational in forecasting future states of the world, however, the households are rational in solving their own maximisation problem. We do this to ensure that our model supports the findings of \cite{fagereng_mpc_2021} who find that households have different marginal propensities to consume depending on their level of income and assets. Each period the households update their expectation beliefs and consequently recalculate their optimal choices given their updated beliefs. Banks play the role of financial intermediaries by financing firms and households. The central bank sets the interest rate according to a \cite{taylor_discretion_1993} rule. The government levies taxes and distributes social support and buys government consumption. \\

\noindent Typical ingredients of agent based models are heterogeneity, bounded rationality, disequilibrium, and non-linearities \citep{axtell_agent-based_nodate}. Agents' behaviour are specified by simple rules which are conditional on the state of the agent and global states although these rules can be generalized to more complex or even optimal ones. The dependency of the agents' behaviour on the agents' states results in heterogeneous behaviour. The agents are embedded in an external environment such as markets or networks which determines the interactions of the agents. Aggregate states are not pre-specified but emerge from the interaction of the agents. 

Large macroeconomic agent based models are a more recent development. For instance,\cite{riccetti_agent_2015} developed a macroeconomic agent based model and found that their model generated endogenous dynamics such as business cycles, endogenous growth, unemployment rate fluctuations, leverage cycles, bank defaults, and financial instability. Stability transitions endogenously to unstable dynamics and vice versa similar to the \cite{minsky_financial_1999} financial instability hypothesis. \cite{assenza_emergent_2015} created an agent based model with capital and credit and consumption and capital firms. Subsequently, \cite{caiani_agent_2016} developed a stock-flow consistent macroeconomic agent based model which introduced credit networks. More macroeconomic agent based models have been developed, see \cite{dawid_agent-based_2018} for an overview. The previous generation of the macroeconomic agent based models were qualitative as they focused on generating empirical regularities. Current generation models such as the model of \cite{poledna_economic_2019} are quantitative as they have similar or even better forecasting capabilities than traditional methods.\\

\noindent The remainder of this paper is structured as follows. Section 2 gives a short summary of related literature. Section 3 describes the agent based model with extended household sector. Section 4 describes the calibration strategy. Section 5 gives a comparison of the forecast performance of the optimal agent ABM, the base ABM, an extended heuristic ABM, and a VAR model in different forecasting setups. Section 6 presents the analysis of the internal validity by means of the comparison of cross- and auto-correlations of the model to their empirical counterparts. Section 7 concludes the paper.